\subsection{Common Floating-Point Semantics}

The S and D formats are IEEE-754 binary32 and binary64. S uses sign bit 31, exponent bits 30..23, and fraction bits
22..0. D uses sign bit 63, exponent bits 62..52, and fraction bits 51..0. An all-ones exponent and nonzero fraction
is a NaN. The most significant fraction bit distinguishes a quiet NaN from a signaling NaN: zero is signaling and
one is quiet. The architectural default quiet NaNs are \texttt{0x7fc00000} for S and
\texttt{0x7ff8000000000000} for D. Figure~\ref{fig:floating-point-register-encodings} shows both encodings in their
64-bit Fn register view.

An S operation reads bits 31..0 of each Fn operand. Writing an S result to Fn writes the result to bits 31..0 and
writes zero to bits 63..32. A D operation reads and writes all 64 bits. Operations use the selected format for their
intermediate and final values. FMADD, FMSUB, FNMADD, and FNMSUB form the complete multiply-add expression with
unbounded intermediate range and precision and round only the final result. Other arithmetic operations round once
after computing their exact mathematical result.

\subsubsection{Input, NaN, Rounding, and Result Order}

Each ordinary floating-point operation applies the following order.

\begin{enumerate}
\item Read all operands and classify zero, subnormal, normal, infinity, quiet NaN, and signaling NaN.
\item If \texttt{FSTATUS.DAZ} is set, replace each subnormal input with a zero having the same sign.
\item Determine signaling-NaN and operation-specific floating-point exception causes and form the exact mathematical result.
\item For a NaN result, select the first signaling NaN in source-operand order, or the first quiet NaN when there is
no signaling NaN, and set its quiet bit. If the operation produces a NaN without a NaN operand, use the default
quiet NaN. When \texttt{FSTATUS.DN} is set, replace the selected NaN with the default quiet NaN.
\item Round a numeric result to the selected format using \texttt{FSTATUS.RM}, except where an instruction names a
fixed rounding direction.
\item Detect overflow and tininess after rounding. UF is generated when the rounded result is tiny and inexact.
\item If \texttt{FSTATUS.FTZ} is set and the rounded result is a nonzero subnormal, replace it with a zero of the
same sign and generate both UF and NX.
\item Test all generated causes against their FSTATUS enable bits and then perform the architectural commit.
\end{enumerate}

A signaling NaN generates NV. A quiet NaN does not generate NV unless the operation itself is invalid. When two
NaN operands have the same priority, source-operand order determines the selected sign and payload. Quieting sets
the quiet bit and preserves the remaining selected sign and payload bits.

Overflow generates OF and NX. Nearest-even produces signed infinity. Toward zero produces the signed maximum finite
value. Toward positive produces positive infinity for a positive result and the negative maximum finite value for a
negative result. Toward negative produces negative infinity for a negative result and the positive maximum finite
value for a positive result.

\subsubsection{Floating-Point Exception Causes and Commit}

\begin{manualtabular}{@{}>{\raggedright\arraybackslash}p{0.55in}>{\raggedright\arraybackslash}p{4.85in}@{}}
\toprule
\rowcolor{ManualHeaderFill}
\textbf{Cause} & \textbf{Generated when}\\
\midrule
\texttt{NV} & an input is a signaling NaN; an arithmetic operation has no defined numeric result, including infinity minus the same infinity, zero times infinity, zero divided by zero, infinity divided by infinity, or square root of a negative nonzero value; or a floating-to-integer conversion is invalid\\
\texttt{DZ} & a finite nonzero value is divided by zero\\
\texttt{OF} & the rounded numeric magnitude exceeds the selected format's finite range\\
\texttt{UF} & the result is tiny after rounding and inexact, or FTZ flushes a nonzero subnormal result\\
\texttt{NX} & rounding or a valid conversion changes the exact value, or OF or FTZ generates NX\\
\bottomrule
\end{manualtabular}

If any generated floating-point exception cause has its matching FSTATUS enable bit set, the instruction raises the
\texttt{FLOATING\_POINT\_FAULT} architectural event. Its error-code cause bitmap contains every cause generated by the operation.
The destination, integer FLAGS, and FFLAGS remain unchanged, and arithmetic instructions never modify FSTATUS.

If no generated floating-point exception cause is enabled, the instruction commits its destination and any complete
integer FLAGS image defined by the instruction, and ORs every generated cause into FFLAGS. FFLAGS fields not named by
the instruction remain unchanged. An instruction-specific rule may exclude a cause, as the FPTRANSA contracts exclude
NX.

\subsubsection{Comparison, Minimum, and Maximum}

FCMP and FTEST write \texttt{Z,N,C,V} as follows.

\begin{manualtabular}{@{}>{\raggedright\arraybackslash}p{1.35in}>{\centering\arraybackslash}p{0.55in}>{\centering\arraybackslash}p{0.55in}>{\centering\arraybackslash}p{0.55in}>{\centering\arraybackslash}p{0.55in}@{}}
\toprule
\rowcolor{ManualHeaderFill}
\textbf{Relation} & \textbf{Z} & \textbf{N} & \textbf{C} & \textbf{V}\\
\midrule
greater & 0 & 0 & 0 & 0\\
equal & 1 & 0 & 0 & 0\\
less & 0 & 1 & 1 & 0\\
unordered & 0 & 0 & 0 & 1\\
\bottomrule
\end{manualtabular}

A quiet NaN produces unordered without NV. A signaling NaN produces unordered and NV. FTEST compares its operand
with positive zero under these rules.

For FMIN and FMAX, when exactly one operand is a NaN the numeric operand is the result. When both operands are NaNs,
the common NaN-selection rule supplies the result. A signaling NaN generates NV even when the other operand is
numeric. Between positive and negative zero, FMIN selects negative zero and FMAX selects positive zero.

\subsubsection{Conversions}

For Fn-to-Fn FCVT or FCVTU, the \texttt{.S} encoding converts D to S and the \texttt{.D} encoding converts S to D. The two
mnemonics have identical behavior for this conversion family. Rn-to-Fn FCVT interprets all 64 source bits as a
signed integer; FCVTU interprets them as an unsigned integer. Integer-to-floating conversion uses
\texttt{FSTATUS.RM}.

Fn-to-Rn conversion truncates toward zero. FCVT produces a signed 64-bit result and FCVTU produces an unsigned
64-bit result. A valid conversion that discards a fractional part generates NX. An invalid conversion generates NV
without OF or NX. When NV is not enabled, the committed invalid result is:

\begin{manualtabular}{@{}>{\raggedright\arraybackslash}p{1.15in}>{\raggedright\arraybackslash}p{2.05in}>{\raggedright\arraybackslash}p{2.20in}@{}}
\toprule
\rowcolor{ManualHeaderFill}
\textbf{Instruction} & \textbf{Input} & \textbf{Result}\\
\midrule
\texttt{FCVT} & negative overflow or negative infinity & \texttt{0x8000000000000000}\\
\texttt{FCVT} & positive overflow or positive infinity & \texttt{0x7fffffffffffffff}\\
\texttt{FCVT} & NaN & \texttt{0x0000000000000000}\\
\texttt{FCVTU} & negative value or negative infinity & \texttt{0x0000000000000000}\\
\texttt{FCVTU} & positive overflow or positive infinity & \texttt{0xffffffffffffffff}\\
\texttt{FCVTU} & NaN & \texttt{0x0000000000000000}\\
\bottomrule
\end{manualtabular}
